\documentclass[10pt,a4paper]{article}
\usepackage[margin=16mm,top=15mm,bottom=15mm]{geometry}
\usepackage[T1]{fontenc}
\usepackage{lmodern}
\usepackage{microtype}
\usepackage{xcolor}
\usepackage{array,booktabs,tabularx}
\usepackage{fancyhdr}
\usepackage[hidelinks]{hyperref}
\hypersetup{
  pdftitle={Maus Crepe Insights: Project Report},
  pdfsubject={BaselHack 2026 project summary}
}
\usepackage{parskip}
\definecolor{ink}{HTML}{263238}
\definecolor{muted}{HTML}{68777D}
\definecolor{teal}{HTML}{087E78}
\definecolor{tealLight}{HTML}{E7F3F1}
\definecolor{coral}{HTML}{D96B52}
\definecolor{coralLight}{HTML}{FBEEE9}
\definecolor{gold}{HTML}{D49A37}
\definecolor{paper}{HTML}{F5F2EA}
\definecolor{line}{HTML}{D9E0DE}
\color{ink}
\setlength{\parindent}{0pt}
\setlength{\parskip}{4pt}
\setlength{\leftmargini}{14pt}
\setlength{\itemsep}{2pt}
\setlength{\topsep}{2pt}
\makeatletter
\renewcommand\section{\@startsection{section}{1}{0pt}{9pt}{5pt}%
  {\normalfont\large\bfseries\color{ink}}}
\makeatother
\pagestyle{fancy}
\fancyhf{}
\renewcommand{\headrulewidth}{0pt}
\fancyfoot[L]{\footnotesize\color{muted}MAUS CR\^EPE INSIGHTS
  \quad / \quad BASELHACK 2026}
\fancyfoot[R]{\footnotesize\color{muted}\thepage}
\newcommand{\metric}[3]{%
  \colorbox{tealLight}{\begin{minipage}[t][17mm][c]{0.222\linewidth}
    \centering{\color{teal}\bfseries\LARGE #1}\par\vspace{1pt}
    {\color{muted}\scriptsize #2}\par{\color{muted}\tiny #3}
  \end{minipage}}}
\newcommand{\tagline}[1]{{\color{teal}\small\bfseries\MakeUppercase{#1}}}
\newcommand{\source}[1]{{\color{muted}\scriptsize Source: #1\par}}
\begin{document}
\thispagestyle{fancy}

{\color{teal}\rule{\linewidth}{3pt}}\vspace{6pt}
\tagline{Project report \quad / \quad 01 November 2026}\par
\vspace{4pt}
{\fontsize{26}{29}\selectfont\bfseries From run logs to\par
cr\^epe insights}\par
\vspace{4pt}
{\large\color{muted}Maus Cr\^epe Insights \quad | \quad BaselHack 2026}\par
\vspace{8pt}
\colorbox{paper}{\parbox{\dimexpr\linewidth-2\fboxsep}{\small
\textbf{Abstract.} We built an offline-reproducible pipeline to turn Maus
Robotics run records and camera checks into analysis tables, vision
measurements, model artifacts and review tools. The public repository contains
working components, but not a validated end-to-end robot system. Sparse check
coverage and weak labels limit what the reported model scores establish.}}
\vspace{8pt}

\noindent\metric{828}{run records}{exported as JSON and MCAP}%
\hfill\metric{177}{table columns}{one row per run}%
\hfill\metric{138}{vision rows}{of 828 exported runs}%
\hfill\metric{62}{human mask labels}{valid comparisons}

\section{Problem and approach}
Maus robots make cr\^epe-quality decisions from a camera above a hotplate.
Steam, glare, browning, a rotating plate and tools crossing the view make
single-frame inspection difficult. The challenge asks for spread quality,
position and orientation, and fold state, with a path toward real-time use.

\vspace{2pt}
\colorbox{paper}{\parbox{\dimexpr\linewidth-2\fboxsep}{\centering\small\bfseries
Dashboard $\rightarrow$ JSON / MCAP $\rightarrow$ Metrics + models
$\rightarrow$ Review tools}}
\vspace{3pt}

The exporter writes one ROS 2 MCAP per run and a 177-column Parquet table.
The vision pipeline reads the robot's annotated check images as measurements
such as diameter, centre offset, holes and fold overlap. Shared frame
preprocessing feeds compact ONNX classifiers. Separate labeling and mask
prefill tools support future human review; an overlay renderer and run viewer
make outputs inspectable. Private run data, videos and annotations remain out
of Git.

\section{What the measurements say}
\begin{tabularx}{\linewidth}{@{}
  >{\raggedright\arraybackslash\bfseries}p{0.22\linewidth}
  X
  >{\raggedright\arraybackslash}p{0.28\linewidth}@{}}
\toprule
Evidence & Result & Interpretation \\
\midrule
Coverage & 138 of 828 runs have extracted vision measurements; 160 have any
recorded check. & Coverage is selective, not a representative sample of all
runs. Most failed runs lack a recorded check. \\
Segmentation student & Pixel IoU 0.966 against SAM 3 teacher masks on held-out
runs. & Agreement with pseudo-labels, not a human-verified quality score. \\
Human comparison & Median mask IoU 0.944 on 62 valid masks; human--human median
0.933 on 10 matched frames. & Five consecutive runs from one robot; four were
used for student training. This is a pilot, not an independent test. \\
Fold-state rule & 99.5\% agreement at 991 passed-check moments, versus 88.5\%
for the classical pipeline on the same frames. & Thresholds were read from
these frames, including the test subset; results are optimistic. \\
Runtime & 5.71 ms median preprocessing plus inference on a four-thread Intel
CPU. & Raspberry Pi 5 estimate is 64 ms, not a measurement on the board. \\
\bottomrule
\end{tabularx}
\vspace{2pt}
\source{Guide: \texttt{crepe-student.md}; model card: \texttt{model.json};
coverage: \texttt{status.md}, snapshot d1c3701.}
\vspace{7pt}
\colorbox{coralLight}{\parbox{\dimexpr\linewidth-2\fboxsep}{\small
\textbf{Readout.} The project has a reproducible data and model toolchain and a
promising offline segmentation prototype. Its human-label comparison is a
small pilot, not an independent test; the on-robot decision path is not yet
integrated.}}

\newpage
\thispagestyle{fancy}
{\color{coral}\rule{\linewidth}{3pt}}\vspace{6pt}
\tagline{Findings \quad / \quad Limits before deployment}\par
\vspace{4pt}
{\fontsize{22}{25}\selectfont\bfseries A useful prototype,\par
not yet a robot-ready system}\par
\vspace{8pt}

\section{Model results, with the right baseline}
The strongest new result is a compact LR-ASPP MobileNetV3-Large segmentation
student distilled from SAM 3 masks. It produces a cr\^epe mask; fold state is
then estimated from mask area relative to that run's flat-cr\^epe reference.
The exported ONNX model matches its PyTorch mask on the parity sample and is
small enough for CPU inference. On the held-out runs, mask IoU against the
teacher is 0.966. Check-moment agreement is 99.5\% on 991 passed checks, but
that figure is not an independent estimate: thresholds were selected using the
same check moments, including the evaluation subset.

A first comparison against reviewed human labels reports median mask IoU 0.944
on 62 valid masks; two human label sets reach median IoU 0.933 on 10 matched
frames. The label export covers five consecutive runs of one robot, and four
runs also appear in student training. These numbers are a useful review pilot,
not an independent estimate of performance on new robots or runs.

Earlier per-check classifiers predict the robot's own pass/fail verdicts, not
verified food quality. Their holdout balanced accuracy ranges from 0.67 to
0.92 on only 31--61 frames per check; three of four cards say they do not
generalize. These results are weak evidence for deployment and must not be
read as human-validated performance.

\section{Known gaps}
\begin{itemize}
  \item \textbf{Coverage:} vision measurements cover 138/828 runs, all from a
  short recent window in the documented export. Missing measurements are
  unknown, not evidence of a good or bad run.
  \item \textbf{Human evaluation:} the first comparison covers five
  consecutive runs of one robot; four runs were used to train the student.
  Broader, independent evaluation across runs and robots is still needed.
  \item \textbf{Scope:} the student has no hole class, does not estimate
  position or angle, and misses small fresh batter, hands and tools. It was
  trained on check-moment frames.
  \item \textbf{Generalization:} thresholds include evaluation frames; no
  leave-one-robot-out student evaluation has been run, and one robot has only
  six training runs.
  \item \textbf{Integration:} the repository has a mask model, review tools and
  overlay code, but no verified live path that returns all challenge decisions
  from one camera frame. Raspberry Pi performance is estimated, not measured.
\end{itemize}

\section{What is reproducible}
The code, model card, ONNX student, synthetic offline tests and documentation
are public. The private dashboard export, frame tables, teacher masks and
labeling records are not in the repository, so a fresh clone can reproduce the
code checks but not the data-dependent counts or training run.

\colorbox{tealLight}{\parbox{\dimexpr\linewidth-2\fboxsep}{\small
\textbf{Offline check}\par
\texttt{cd code \&\& pixi install \&\& pixi run test \&\& pixi run lint}\par
This checks the public code; it does not rebuild private data artifacts.}}

\section{Evidence trail}
\small
\begin{tabularx}{\linewidth}{@{}
  >{\raggedright\arraybackslash\bfseries}p{0.34\linewidth}X@{}}
\toprule
Claim & Repository source \\
\midrule
Data counts and known gaps & \texttt{documentation/status.md} (snapshot
revision d1c3701; counts from private export) \\
Student method and evaluation &
\texttt{documentation/vision/crepe-student.md} \\
Model shape, splits and metrics &
\texttt{code/models/crepe-student/model.json} \\
Model use and rebuild requirements &
\texttt{code/models/crepe-student/README.md} \\
Project source &
\href{https://github.com/leonardoAB1/maus-crepe-insights}
{project repository} \\
\bottomrule
\end{tabularx}

\vfill
{\color{line}\rule{\linewidth}{0.6pt}}\par
{\scriptsize\color{muted}Evidence snapshot: public repository main at
\texttt{f5d1504} (1 Nov 2026). Data coverage counts are explicitly tied to the
older status snapshot because that page has not been refreshed for this
revision. The model-specific results above come from the student guide and
model card at the evidence snapshot. This report is a concise account, not a
claim of production readiness.}
\end{document}
